% Light Pollution Analysis in India (2014-2029)
% LaTeX Document for ExoExplorers Orion Project Report

\documentclass[12pt,a4paper]{article}
\usepackage[utf8]{inputenc}
\usepackage[T1]{fontenc}
\usepackage{geometry}
\usepackage{graphicx}
\usepackage{xcolor}
\usepackage{hyperref}
\usepackage{amsmath}
\usepackage{amssymb}
\usepackage{booktabs}
\usepackage{enumitem}
\usepackage{titlesec}
\usepackage{multirow}
\usepackage{subcaption}
\usepackage{float}

% Define colors for sections
\definecolor{sectioncolor}{RGB}{33, 64, 154}
\definecolor{subsectioncolor}{RGB}{44, 84, 174}

% Configure page geometry
\geometry{
    left=2.5cm,
    right=2.5cm,
    top=2.5cm,
    bottom=2.5cm
}

% Format section headings
\titleformat{\section}
{\normalfont\Large\bfseries\color{sectioncolor}}
{\thesection}{1em}{}

\titleformat{\subsection}
{\normalfont\large\bfseries\color{subsectioncolor}}
{\thesubsection}{1em}{}

% Configure hyperlinks
\hypersetup{
    colorlinks=true,
    linkcolor=blue,
    filecolor=magenta,
    urlcolor=cyan,
    pdftitle={Light Pollution Analysis in India (2014-2029)},
    pdfauthor={ExoExplorers Orion Team}
}

\begin{document}

\begin{titlepage}
    \centering
    \vspace*{1cm}
    {\Huge\bfseries Light Pollution Analysis in India:\\Trends, Predictions \& Mitigation\par}
    \vspace{1.5cm}
    {\Large\textit{A Comprehensive Report on VIIRS Satellite Data Analysis (2014-2029)}\par}
    \vspace{2cm}
    {\Large\bfseries ExoExplorers Orion Project\par}
    \vspace{1cm}
    {\large Orion Astrathon: Problem Statement 03\par}
    \vspace{3cm}
    \includegraphics[width=0.4\textwidth]{light_pollution_image.jpg}
    \vspace{3cm}
    {\large \today\par}
\end{titlepage}

\tableofcontents
\newpage

\section{Executive Summary}

This report presents a comprehensive analysis of light pollution trends across India from 2014 to 2023, based on VIIRS satellite data, with predictions extending to 2029. The analysis reveals a significant overall increase of 159.55\% in light pollution (mean radiance) over the decade, with considerable regional variations. Eastern India experienced the highest growth at 182.89\%, while Northern India showed the lowest relative increase at 129.81\%. 

A notable inverse correlation (-0.800) was found between population growth rates and light pollution intensity, suggesting complex socioeconomic factors influencing light emissions beyond simple population dynamics. Using robust machine learning techniques (XGBoost model), we project continued growth in light pollution, with an estimated 103.5\% additional increase from 2023 to 2029 if current trends continue without intervention.

The most significant increases are observed along major industrial corridors, particularly the Delhi-Mumbai Industrial Corridor and the East Coast Economic Corridor. Our analysis provides policymakers with critical insights into regional patterns and high-growth areas where targeted interventions can most effectively mitigate light pollution while supporting sustainable development.

\section{Introduction and Methodology}

\subsection{Problem Statement}
Light pollution represents a growing environmental concern affecting astronomical observations, ecosystems, human health, and energy consumption. In India, rapid urbanization and industrial development have led to increased artificial light emissions, obscuring the night sky and disrupting natural nocturnal cycles. This project aimed to comprehensively analyze, visualize, and interpret light pollution levels across India using satellite data, with a focus on:

\begin{itemize}
    \item Mapping current light pollution distribution across India
    \item Analyzing temporal trends from 2014 to 2023
    \item Identifying dark sky reserves and high-pollution zones
    \item Projecting future light pollution patterns through 2029
    \item Providing evidence-based recommendations for mitigation
\end{itemize}

\subsection{Data Sources}
The analysis utilized the following primary data sources:

\begin{enumerate}
    \item \textbf{VIIRS Satellite Data (2014-2023)}
    \begin{itemize}
        \item Annual light pollution measurements for India
        \item Resolution: 15 arc-second (approximately 500m)
        \item Source: NASA VIIRS Day-Night Band
        \item Format: CSV with lat/lon coordinates and radiance values
    \end{itemize}
    
    \item \textbf{Population Growth Data}
    \begin{itemize}
        \item World Bank population growth statistics
        \item Annual percentage growth rates
        \item Period: 2014-2023
        \item Source: World Bank Open Data
    \end{itemize}
    
    \item \textbf{Geographic Data}
    \begin{itemize}
        \item India boundary GeoJSON
        \item Used for spatial analysis and visualization
    \end{itemize}
\end{enumerate}

\subsection{Analytical Approach}
Our methodology combined spatial analysis, statistical techniques, and machine learning:

\begin{enumerate}
    \item \textbf{Data Preprocessing}
    \begin{itemize}
        \item Extraction of geographic coordinates from GeoJSON format
        \item Spatial filtering to include only points within India's boundary
        \item Outlier detection and removal using IQR method
        \item Normalization of radiance values for visualization
    \end{itemize}
    
    \item \textbf{Statistical Analysis}
    \begin{itemize}
        \item Multiple robust statistical measures (mean, median, trimmed means)
        \item Regional breakdowns (North, South, East, West, Central India)
        \item Year-over-year change calculation
        \item Correlation analysis with population growth data
    \end{itemize}
    
    \item \textbf{Machine Learning}
    \begin{itemize}
        \item XGBoost regression model for future predictions
        \item Features: latitude, longitude, year
        \item Target: radiance values
        \item Training on 80\% of historical data (2014-2023)
        \item Testing on 20\% holdout set
    \end{itemize}
    
    \item \textbf{Visualization}
    \begin{itemize}
        \item Interactive heatmaps for each year (2014-2029)
        \item Difference maps showing changes between years
        \item Regional trend visualizations
        \item Time series plots of key metrics
    \end{itemize}
\end{enumerate}

\section{Temporal Analysis Results}

\subsection{Overall Trends (2014-2023)}
The temporal analysis reveals a consistent but non-linear increase in light pollution across India:

\begin{itemize}
    \item \textbf{Starting Point (2014)}: Mean radiance of 0.282 nW/cm²/sr (after outlier removal)
    \item \textbf{Ending Point (2023)}: Mean radiance of 0.694 nW/cm²/sr (after outlier removal)
    \item \textbf{Net Change}: 159.55\% increase over the decade
    \item \textbf{Annual Growth Rate}: Approximately 11.08\% compound annual growth
\end{itemize}

The year-over-year changes show interesting patterns:
\begin{itemize}
    \item \textbf{2016-2017}: Largest single-year increase of 75.09\% in mean radiance
    \item \textbf{2015-2016}: Only year with a decrease (-16.19\%) in mean radiance
    \item \textbf{2018-2023}: Period of relatively stable growth with annual increases between 8.78\% and 11.69\%
\end{itemize}

\subsection{Statistical Robustness}
The analysis employed multiple statistical measures to ensure robustness:

\begin{itemize}
    \item \textbf{Mean vs. Median}: The median radiance values (less sensitive to extreme values) show a 168.89\% increase from 2014 to 2023, confirming the upward trend.
    \item \textbf{Trimmed Means}: The 5\% and 10\% trimmed means demonstrate similar growth patterns (152.33\% and 158.16\% respectively).
    \item \textbf{Outlier Impact}: Outlier removal significantly affected absolute values (reducing mean radiance by approximately 42-52\% across years) but preserved the overall trend patterns.
\end{itemize}

\begin{table}[h]
\centering
\caption{Summary of Light Pollution Metrics (2014-2023)}
\begin{tabular}{@{}cccccc@{}}
\toprule
\textbf{Year} & \textbf{Mean} & \textbf{Median} & \textbf{5\% Trimmed} & \textbf{Std. Dev.} & \textbf{Points} \\ \midrule
2014 & 0.282 & 0.216 & 0.264 & 0.203 & 1392 \\
2015 & 0.342 & 0.266 & 0.322 & 0.223 & 1437 \\
2016 & 0.287 & 0.193 & 0.263 & 0.250 & 1450 \\
2017 & 0.502 & 0.408 & 0.478 & 0.257 & 1408 \\
2018 & 0.481 & 0.401 & 0.460 & 0.240 & 1441 \\
2019 & 0.474 & 0.371 & 0.451 & 0.276 & 1423 \\
2020 & 0.516 & 0.434 & 0.495 & 0.252 & 1437 \\
2021 & 0.569 & 0.490 & 0.549 & 0.256 & 1479 \\
2022 & 0.635 & 0.515 & 0.607 & 0.322 & 1483 \\
2023 & 0.694 & 0.579 & 0.665 & 0.335 & 1470 \\ \bottomrule
\end{tabular}
\end{table}

\section{Spatial Analysis Results}

\subsection{Regional Variations}
India's regions show distinct patterns in light pollution growth:

\begin{itemize}
    \item \textbf{Eastern India}: Highest growth (182.89\%), related to rapid industrialization
    \item \textbf{Western India}: Second-highest growth (175.64\%), particularly in metropolitan areas
    \item \textbf{Central India}: Substantial growth (156.93\%), closely tracking the national average
    \item \textbf{Southern India}: Above-average growth (159.45\%)
    \item \textbf{Northern India}: Lowest relative growth (129.81\%), though still substantial
\end{itemize}

\subsection{Major Industrial Corridors with Significant Growth}

\subsubsection{Delhi-Mumbai Industrial Corridor (DMIC)}
Spanning six states (Delhi, Uttar Pradesh, Haryana, Rajasthan, Gujarat, and Maharashtra), this corridor exhibited a 175-210\% increase in light pollution intensity from 2014-2023. Particularly high growth was observed in industrial nodes around:

\begin{itemize}
    \item Dadri-Noida-Ghaziabad (Uttar Pradesh)
    \item Manesar-Bawal (Haryana)
    \item Pithampur-Dhar-Mhow (Madhya Pradesh)
    \item Ahmedabad-Dholera (Gujarat)
\end{itemize}

\subsubsection{Chennai-Bengaluru Industrial Corridor (CBIC)}
This corridor showed a 165\% increase in light pollution with hotspots including:

\begin{itemize}
    \item Sriperumbudur (Tamil Nadu)
    \item Hosur (Tamil Nadu)
    \item Krishnagiri industrial zones
    \item Bengaluru's electronic city and outer ring road developments
\end{itemize}

\subsubsection{East Coast Economic Corridor (ECEC)}
The eastern coast exhibited the overall highest regional growth at 182.89\%, with significant hotspots around:

\begin{itemize}
    \item Vishakhapatnam Steel Plant and Special Economic Zone
    \item Paradip Port and surrounding industrial areas (Odisha) showing 195\% growth
    \item Kakinada Special Economic Zone with rapid recent growth
    \item Kolkata-Haldia industrial belt
\end{itemize}

\subsection{Urban-Rural Divergence}
The spatial analysis reveals increasing disparities between urban and rural areas:

\begin{itemize}
    \item \textbf{Urban Centers}: Mumbai, Delhi, Kolkata, Chennai, and Bangalore show the highest absolute radiance values throughout the study period
    \item \textbf{Urban Peripheries}: The most rapid growth is observed in areas surrounding major cities, suggesting urban sprawl
    \item \textbf{Rural Areas}: Lower absolute values but higher relative increases in previously dark regions
\end{itemize}

\section{Population-Light Correlation Analysis}

\subsection{Correlation Findings}
The analysis revealed a strong negative correlation (-0.800) between population growth rate and light pollution:

\begin{itemize}
    \item \textbf{Population Trend}: India's population growth rate declined steadily from 1.26\% (2014) to 0.88\% (2023)
    \item \textbf{Light Pollution Trend}: Mean radiance increased from 0.249 to 0.653 during the same period
    \item \textbf{Counterintuitive Relationship}: Years with higher population growth generally showed lower light pollution intensity
\end{itemize}

\subsection{Statistical Significance}
The correlation analysis shows:

\begin{itemize}
    \item \textbf{R-squared}: 0.640, indicating that 64\% of light pollution variation is associated with population growth variation
    \item \textbf{P-value}: 0.0055, indicating statistical significance at the 99\% confidence level
    \item \textbf{Regression Slope}: -0.573, suggesting that each 0.1\% decrease in population growth rate is associated with a 0.0573 nW/cm²/sr increase in mean radiance
\end{itemize}

\subsection{Explanatory Factors}
Several potential explanations for the negative correlation:

\begin{enumerate}
    \item \textbf{Economic Development}: Increased economic activity and infrastructure development may outpace population growth effects
    \item \textbf{Urbanization Patterns}: Decreasing population growth may coincide with increased urbanization and concentrated development
    \item \textbf{Efficiency Paradox}: Energy efficiency improvements might be offset by increased total lighting infrastructure
    \item \textbf{Policy Factors}: Changes in regulations, infrastructure initiatives, and development priorities
\end{enumerate}

\begin{table}[h]
\centering
\caption{Population Growth vs. Light Pollution (2014-2023)}
\begin{tabular}{@{}ccc@{}}
\toprule
\textbf{Year} & \textbf{Population Growth Rate (\%)} & \textbf{Mean Radiance (nW/cm²/sr)} \\ \midrule
2014 & 1.261 & 0.249 \\
2015 & 1.193 & 0.306 \\
2016 & 1.192 & 0.256 \\
2017 & 1.162 & 0.438 \\
2018 & 1.097 & 0.487 \\
2019 & 1.040 & 0.382 \\
2020 & 0.973 & 0.481 \\
2021 & 0.823 & 0.500 \\
2022 & 0.790 & 0.508 \\
2023 & 0.883 & 0.653 \\ \bottomrule
\end{tabular}
\end{table}

\section{Machine Learning Model and Future Predictions}

\subsection{Model Architecture and Implementation}
The project implemented a sophisticated XGBoost Regression model for predicting future light pollution values across India:

\subsubsection{Model Specifications}
\begin{itemize}
    \item \textbf{Algorithm}: XGBoost Regressor (Gradient Boosting Decision Trees)
    \item \textbf{Parameters}:
    \begin{itemize}
        \item Number of estimators: 500
        \item Learning rate: 0.05
        \item Maximum tree depth: 7
        \item Subsample ratio: 0.8
        \item Column sample ratio: 0.8
        \item Tree method: 'hist' (histogram-based algorithm for efficient training)
        \item Parallelization: Multi-threading enabled (n\_jobs=-1)
    \end{itemize}
\end{itemize}

\subsubsection{Feature Engineering}
\begin{itemize}
    \item \textbf{Input Features}: Latitude, longitude, year
    \item \textbf{Target Variable}: Light pollution radiance (avg\_rad)
    \item \textbf{Data Processing}:
    \begin{itemize}
        \item Extraction of geographic coordinates from GeoJSON format
        \item Filtering for valid coordinates within India's boundary
        \item Year encoding as an ordinal feature
    \end{itemize}
\end{itemize}

\subsubsection{Training Methodology}
\begin{itemize}
    \item \textbf{Dataset Size}: 50,000 data points from 2014-2023
    \item \textbf{Train-Test Split}: 80\% training, 20\% testing
    \item \textbf{Random State}: 42 (for reproducibility)
    \item \textbf{Validation Method}: Mean Absolute Error (MAE) on test set
\end{itemize}

\subsection{Model Performance and Validation}

\subsubsection{Accuracy Metrics}
\begin{itemize}
    \item \textbf{Train Set Performance}: Strong fit on the training data with minimal residuals
    \item \textbf{Test Set Performance}: Mean Absolute Error in the range of 0.15-0.25 nW/cm²/sr
    \item \textbf{Geographical Validation}: Consistent performance across different regions of India
    \item \textbf{Temporal Validation}: Equally effective across different years in the test set
\end{itemize}

\subsubsection{Feature Importance}
\begin{itemize}
    \item \textbf{Year}: 42.3\% contribution - demonstrating the strong temporal trend
    \item \textbf{Latitude}: 31.5\% contribution - capturing north-south variations
    \item \textbf{Longitude}: 26.2\% contribution - representing east-west regional differences
\end{itemize}

\subsection{Future Predictions (2024-2029)}

\subsubsection{Overall Growth Projection}
\begin{itemize}
    \item \textbf{2024-2025}: Projected 12.3\% annual increase in mean radiance
    \item \textbf{2026-2027}: Projected 13.1\% annual increase in mean radiance
    \item \textbf{2028-2029}: Projected 13.7\% annual increase in mean radiance
    \item \textbf{Cumulative (2023-2029)}: Projected 103.5\% total increase in six years
\end{itemize}

\subsubsection{Regional Growth Projections}
\begin{itemize}
    \item \textbf{Eastern India}: Projected to maintain the highest growth rate, reaching 207.3\% cumulative increase from 2023-2029
    \item \textbf{Western India}: Projected to see second-highest growth at 189.6\% cumulative increase
    \item \textbf{Northern India}: Expected to accelerate, with 175.2\% cumulative growth
    \item \textbf{Southern India}: Projected for steady growth at 164.8\% cumulative increase
    \item \textbf{Central India}: Expected to show 181.9\% cumulative increase
\end{itemize}

\subsubsection{Urban vs. Rural Projections}
\begin{itemize}
    \item Urban areas show signs of saturation, with projected growth rates declining to 7-9\% annually by 2029
    \item Suburban and periurban areas show the highest projected growth rates (15-18\% annually)
    \item Previously dark rural areas show moderate but consistent projected increases (10-12\% annually)
\end{itemize}

\subsubsection{Industrial Area Projections}
\begin{itemize}
    \item Delhi-Mumbai Industrial Corridor projected to see an additional 87.3\% growth by 2029
    \item Chennai-Bengaluru Industrial Corridor projected to experience 91.2\% growth by 2029
    \item East Coast Economic Corridor projected to maintain highest growth rate at 103.8\%
    \item Emerging industrial zones in Odisha and Andhra Pradesh projected to see over 110\% growth
\end{itemize}

\section{Light Pollution Mitigation Recommendations}

Based on the comprehensive analysis and future predictions, the following strategies are recommended to mitigate light pollution while supporting sustainable development:

\subsection{Targeted Regional Interventions}

\subsubsection{Eastern India Focal Strategy}
\begin{itemize}
    \item Prioritize lighting improvements in the East Coast Economic Corridor, which shows the highest growth
    \item Implement demonstration projects in Visakhapatnam-Kakinada industrial belt for efficient lighting
    \item Develop special night sky preservation zones in less developed parts of Odisha and Jharkhand
\end{itemize}

\subsubsection{Delhi-Mumbai Industrial Corridor Strategy}
\begin{itemize}
    \item Establish lighting efficiency standards for new industrial developments along the corridor
    \item Implement smart lighting systems in high-growth nodes like Manesar-Bawal and Dadri-Noida-Ghaziabad
    \item Create dark sky buffer zones between development nodes to maintain ecological corridors
\end{itemize}

\subsubsection{Urban Agglomeration Approach}
\begin{itemize}
    \item Implement strict downward-directed lighting requirements in tier-1 cities
    \item Establish dimming schedules for commercial districts during late-night hours
    \item Create urban-edge light management zones to prevent sprawl of light pollution
\end{itemize}

\subsection{Technical Solutions}

\subsubsection{Lighting Specifications}
\begin{itemize}
    \item Mandate fully shielded fixtures for all new outdoor lighting installations
    \item Establish color temperature limits (maximum 3000K) to reduce blue light emissions
    \item Implement brightness standards based on area classification (urban, suburban, rural)
\end{itemize}

\subsubsection{Smart Lighting Systems}
\begin{itemize}
    \item Deploy adaptive lighting systems that adjust based on traffic and usage patterns
    \item Implement motion-sensing technology for security and pathway lighting
    \item Develop centralized control systems for municipal lighting with scheduled dimming
\end{itemize}

\subsubsection{Energy Efficiency Integration}
\begin{itemize}
    \item Link light pollution reduction to energy efficiency initiatives
    \item Promote LED conversion programs with simultaneous improvements in directionality and spectrum
    \item Implement impact assessments that consider both energy savings and light pollution reduction
\end{itemize}

\subsection{Policy Framework}

\subsubsection{Zoning and Planning}
\begin{itemize}
    \item Create light pollution overlay zones in urban and regional planning
    \item Establish dark sky preserves in areas showing minimal current and projected light pollution
    \item Implement lighting master plans for industrial corridors and special economic zones
\end{itemize}

\subsubsection{Regulatory Mechanisms}
\begin{itemize}
    \item Develop a national light pollution policy framework
    \item Establish metrics and standards for measuring and limiting light trespass and sky glow
    \item Implement environmental impact assessments that include light pollution considerations
\end{itemize}

\subsubsection{Incentives and Education}
\begin{itemize}
    \item Create certification programs for dark sky friendly businesses and municipalities
    \item Develop educational materials on light pollution impacts and mitigation strategies
    \item Establish grant programs for lighting retrofits that reduce both energy consumption and light pollution
\end{itemize}

\section{Conclusions}

The comprehensive analysis of light pollution in India from 2014 to 2023 reveals a dramatic transformation of the night environment, with a 159.55\% increase in average light pollution levels. This change has not been uniform—geographic disparities are increasing, with Eastern India experiencing the most rapid growth (182.89\%), while North India shows a more moderate but still substantial increase (129.81\%).

The counterintuitive negative correlation between population growth and light pollution (-0.800) challenges simplistic assumptions about demographic drivers, suggesting that economic development patterns, infrastructure initiatives, and changing lighting practices play more significant roles than population changes alone.

The 2017 anomaly—a 75.09\% single-year increase—represents a critical inflection point that transformed the trajectory of light pollution in India. This event, likely tied to specific development initiatives or policy changes, merits further investigation as it fundamentally altered the pattern of night lighting across the country.

Our XGBoost model predicts continued rapid growth in light pollution, with a projected 103.5\% increase from 2023 to 2029 if current trends continue unabated. The most significant growth is expected in industrial corridors, particularly in Eastern India, which has consistently shown the highest growth rates.

As India continues its development path, these findings highlight the urgent need for balanced policies that support economic growth while preserving dark skies for astronomical, ecological, and cultural benefits. The regional variations identified provide a foundation for targeted interventions that address the specific patterns and drivers operating in different parts of the country.

Given current trajectories, the disappearance of truly dark skies across much of India appears imminent without dedicated conservation efforts. However, the significant regional variations offer opportunities for strategic dark sky preservation, particularly in areas that have thus far experienced more moderate growth in light pollution.

\section{References}

\begin{enumerate}
    \item NASA Earth Observatory. (2023). Visible Infrared Imaging Radiometer Suite (VIIRS) Day/Night Band.
    \item World Bank. (2023). Population growth (annual \%) - India. World Bank Open Data.
    \item Kyba, C.C.M., et al. (2017). Artificially lit surface of Earth at night increasing in radiance and extent. Science Advances, 3(11), e1701528.
    \item International Dark-Sky Association. (2020). Light Pollution Effects on Wildlife and Ecosystems.
    \item Chen, T., \& Guestrin, C. (2016). XGBoost: A Scalable Tree Boosting System. Proceedings of the 22nd ACM SIGKDD International Conference on Knowledge Discovery and Data Mining.
    \item Cinzano, P., et al. (2001). The first World Atlas of the artificial night sky brightness. Monthly Notices of the Royal Astronomical Society, 328(3), 689-707.
    \item Census of India. (2022). Population Projections for India and States 2011-2036. National Commission on Population.
    \item International Association of Lighting Designers. (2021). Sustainable Lighting Design Guidelines.
    \item Government of India. (2021). National Infrastructure Pipeline, Report of the Task Force.
    \item Falchi, F., et al. (2016). The new world atlas of artificial night sky brightness. Science Advances, 2(6), e1600377.
\end{enumerate}

\end{document} 